\documentclass[10pt,letterpaper,oneside]{article}
\usepackage[T1]{fontenc}
\usepackage[utf8]{inputenc}
\usepackage{lmodern}
\usepackage{geometry}
\geometry{margin=0.84in,headheight=21pt,footskip=28pt}
\usepackage{microtype}
\usepackage{xcolor}
\definecolor{ink}{HTML}{15253D}
\definecolor{accent}{HTML}{086A84}
\definecolor{pale}{HTML}{EEF5F8}
\definecolor{muted}{HTML}{4E6474}
\definecolor{grayline}{HTML}{D7E2E8}
\usepackage{graphicx}
\usepackage{fancyhdr}
\usepackage{titlesec}
\usepackage{enumitem}
\usepackage{array,tabularx,longtable,booktabs}
\usepackage[most]{tcolorbox}
\usepackage{xurl}
\usepackage{hyperref}
\usepackage{listings}
\usepackage{amsmath,amssymb}
\usepackage{tikz}
\usetikzlibrary{arrows.meta,positioning,shapes.geometric,calc}
\hypersetup{colorlinks=true,linkcolor=accent,urlcolor=accent,pdftitle={Network, Host and Cloud Defense Operations},pdfauthor={Independent technical engineering study companion}}
\renewcommand{\familydefault}{\sfdefault}
\linespread{1.07}
\setlength{\parindent}{0pt}
\setlength{\parskip}{5.5pt}
\emergencystretch=2em
\setlist[itemize]{leftmargin=1.5em,itemsep=2.5pt,topsep=3pt}
\setlist[enumerate]{leftmargin=1.7em,itemsep=3pt,topsep=3pt}
\titleformat{\section}{\Large\bfseries\color{ink}}{\thesection}{0.65em}{}[\vspace{-2pt}{\color{accent}\titlerule}]
\titleformat{\subsection}{\large\bfseries\color{accent}}{\thesubsection}{0.6em}{}
\titleformat{\subsubsection}{\bfseries\color{ink}}{\thesubsubsection}{0.6em}{}
\titlespacing*{\section}{0pt}{16pt}{9pt}
\titlespacing*{\subsection}{0pt}{11pt}{6pt}
\pagestyle{fancy}\fancyhf{}\fancyhead[L]{\footnotesize\color{muted} COMPUTER \& INFORMATION SECURITY ENGINEERING}\fancyhead[R]{\footnotesize\color{accent} SUITE / 4}\fancyfoot[L]{\footnotesize\color{muted} Independent study companion \quad|\quad 2026}\fancyfoot[R]{\footnotesize\color{muted}\thepage}\renewcommand{\headrulewidth}{0.4pt}\renewcommand{\headrule}{\hbox to\headwidth{\color{grayline}\leaders\hrule height \headrulewidth\hfill}}
\newtcolorbox{sourcebox}{enhanced,breakable,boxrule=0pt,leftrule=3pt,colframe=accent,colback=pale,arc=0pt,left=9pt,right=7pt,top=7pt,bottom=7pt}
\newtcolorbox{warningbox}{enhanced,breakable,boxrule=0.5pt,colframe=grayline,colback=white,arc=1pt,left=9pt,right=9pt,top=7pt,bottom=7pt}
\newcommand{\handbookref}[1]{\textsf{\color{accent}\textbf{[H:#1]}}}
\newcommand{\modernext}{\textsf{\color{muted}\textbf{[E]}}}
\newcommand{\designed}{\textsf{\color{accent}\textbf{[D]}}}
\newcommand{\deliverable}[1]{\textbf{Deliverable:} #1}
\lstdefinestyle{lab}{basicstyle=\ttfamily\scriptsize,keywordstyle=\color{accent}\bfseries,commentstyle=\color{muted},numbers=left,numberstyle=\tiny\color{muted},numbersep=7pt,breaklines=true,showstringspaces=false,frame=single,rulecolor=\color{grayline},backgroundcolor=\color{pale},xleftmargin=1em,framexleftmargin=0.6em,tabsize=2,columns=fullflexible}
\lstset{style=lab}
\renewcommand{\arraystretch}{1.21}
\newcommand{\makereportcover}[4]{%
\thispagestyle{empty}\vspace*{1.15in}
{\small\bfseries\color{accent} COMPUTER AND INFORMATION SECURITY ENGINEERING}\par\vspace{13pt}
{\Huge\bfseries\color{ink} #1\par}\vspace{15pt}
{\Large\color{muted} #2\par}\vspace{18pt}
\color{accent}\rule{0.94\textwidth}{2.1pt}\par\vspace{20pt}
{\normalsize\bfseries Document #3 of 8 \quad\textbullet\quad 16 October 2026\par}
\vspace{0.35in}
\begin{sourcebox}\textbf{Orientation.} #4\end{sourcebox}
\vfill
{\small\color{muted} Independent, original engineering study and practice materials grounded in\par
\emph{Computer and Information Security Handbook}, 4th ed., edited by John R. Vacca,\par
Elsevier/Morgan Kaufmann, 2025, Volumes I--II. Not a reproduction or official standard.}
\newpage\setcounter{page}{1}\tableofcontents\newpage
}
\begin{document}
\makereportcover{Network, Host and Cloud Defense Operations}{Telemetry, segmentation, detection engineering, incident handling and resilience}{4}{Turn the handbook's network monitoring, intrusion prevention and practical security coverage into operating procedures with testable detection and containment mechanisms.}

\section{Defense operating model}
The handbook links observation, alert management, prevention, network protection, TCP/IP packet analysis, endpoints, firewalls, cloud and virtualization. \handbookref{5--8,11--18,64--70,74--76,78,80} This volume converts those concepts into a repeatable detect--investigate--respond--improve loop.
\begin{sourcebox}\textbf{Operating loop:} enumerate assets and high-value flows; select telemetry; write detections; validate normal and abnormal cases; triage alerts; scope incidents; authorize containment; verify recovery; update models and rules. \designed{}\end{sourcebox}

\section{Telemetry strategy and chain of observation}
Chapter 5 foregrounds full packet capture (FPC) and network security monitoring (NSM) as sources of forensic visibility; complete capture, however, raises privacy, retention, throughput and encryption-related constraints. \handbookref{5,44} \modernext{} Prefer collection choices tied to explicit investigative questions rather than indiscriminate capture.
\par\medskip\noindent\begin{tabularx}{\textwidth}{@{}p{0.19\textwidth}X X@{}}\toprule
Signal & High-value question & Blind spot or caveat \\ \midrule
Flow logs & Who connected to which destination? & Little/no application payload detail\\
DNS resolution & Which names did a workload resolve? & Encrypted DNS, caching, non-DNS flows\\
HTTP/TLS metadata & Service, status and certificate behavior? & End-to-end encryption, sampling\\
Endpoint events & Which process/user created an action? & Agent tampering and incomplete coverage\\
Cloud audit & Which principal changed infrastructure? & Coverage depends on enabled services\\
Full packet capture & What bytes traversed an observed link? & Sensitive data, packet loss and encryption\\
Identity logs & Was a privilege/token used abnormally? & Session hijack vs legitimate account use\\
\bottomrule\end{tabularx}\par\medskip
For each feed record source, time synchronization, retention, integrity protection, access rights, parsing health, volume and cost. Silence is an operational fault, not confirmation of safety. \designed{}

\section{Segmentation, ingress/egress and firewall governance}
\handbookref{8,13,15,16,59,70,76} A firewall policy should name a specific legitimate communication. Every broad rule needs justification, expiry and compensating observation.
\par\medskip\noindent\begin{tabularx}{\textwidth}{@{}p{0.21\textwidth}p{0.27\textwidth}X@{}}\toprule
Flow & Default posture & Allowed exception and review \\ \midrule
Internet to workload & Deny & Gateway 443 to published service only\\
Workload to DB & Deny & App principal and database port only\\
Workload to internet & Deny & Approved repository/API destinations\\
Admin to production & Deny & Managed jump access with recorded approval\\
Telemetry outbound & Restricted & TLS to trusted collector only\\
Backup management & Isolated & Scheduled, privileged controlled path\\
\bottomrule\end{tabularx}\par\medskip
\subsection{Change test}\designed{}
For every rule change, test permitted paths, forbidden paths, return traffic and logging. Include emergency rollback and out-of-band recovery; never apply restrictive host firewall configuration to an unprotected remote session.

\section{Intrusion detection and alert engineering}
\subsection{Detection specification}\designed{}
\begin{warningbox}\textbf{DET-014 --- anomalous administrative authentication}\\
\textbf{Hypothesis:} compromised privileged account is used from an unauthorized context.\\
\textbf{Inputs:} successful/failed login, asset criticality, identity, device, session history.\\
\textbf{Logic:} high-risk sign-in + privileged action within a short interval.\\
\textbf{Severity:} depends on privilege, evidence of session theft and affected assets.\\
\textbf{Exclusions:} approved emergency maintenance and defined test identities.\\
\textbf{Response:} preserve audit events; validate with owner; revoke/restrict only under authorization.\\
\textbf{Tests:} known-good privileged session; simulated unauthorized location; log missingness.\\
\textbf{Metrics:} precision of reviewed alerts, time-to-triage, false-negative scenario coverage.
\end{warningbox}
Map detections to relevant attack behaviors and available data before claiming ATT\&CK technique coverage; writing a rule does not demonstrate it works. \handbookref{5,6,74} \modernext{}

\section{Host and container defense}
\handbookref{11,12,69,78,80} Prefer hardened images, security updates, protected administrative paths, minimal listening services, immutable deployment and runtime auditing. Security Onion/Sguil examples in the handbook are valuable historical illustrations; treat their exact installation instructions and obsolete package versions as environment-specific, not timeless deployment recommendations. \modernext{}
\subsection{Daily operational assertions}\designed{}
\begin{itemize}
\item In-scope endpoints are reporting with recent heartbeat and defined telemetry quality.
\item Supported systems have recorded patch status, exception owners and deadlines.
\item Unauthorized local administrators and stale tokens are investigated.
\item Public endpoints and outbound dependencies match the approved inventory.
\item Security instrumentation cannot silently disable itself without alerting.
\end{itemize}

\section{Cloud operations and compromise paths}
\handbookref{9,64--70} Cloud incidents frequently involve identities, configuration, storage exposure or abused APIs as much as compromised hosts. Build a response playbook for a leaked deployment key: verify the key, record its scope, identify access events, restrict/rotate credential, inspect affected resources, restore trusted automation and verify no further unauthorized use. \designed{}

\section{Incident triage: evidence before irreversible action}
\begin{enumerate}
\item \textbf{Intake:} source, timestamp, detection ID, affected asset, confidence and reporter.
\item \textbf{Validate:} confirm event mechanics, identify alternative explanations and check telemetry completeness.
\item \textbf{Scope:} identities, hosts, services, data objects, timelines and likely lateral access.
\item \textbf{Preserve:} logs, volatile data and system state with access control and custody metadata.
\item \textbf{Contain:} select proportionate reversible actions; seek incident command authorization.
\item \textbf{Eradicate and recover:} remove root cause, rotate trust material, redeploy, retest.
\item \textbf{Learn:} update detections, preventive controls, metrics and playbooks.
\end{enumerate}
NIST SP 800-61 Rev.~3 places incident response within the broader CSF 2.0 risk management functions rather than only a linear incident phase chart. \modernext{} [N3]

\section{Sector-specific detection and resilience considerations}
\par\medskip\noindent\begin{tabularx}{\textwidth}{@{}p{0.25\textwidth}X@{}}\toprule
Environment & Required adjustment \\ \midrule
Smart-home/sensor network & Constrained battery/CPU; wireless jamming; device identity and physical access.\\
Energy and grid & Safety and availability outrank intrusive scanning; coordinated change windows.\\
Connected vehicle & Treat external telematics and actuation as separate trust domains.\\
Maritime/satellite & Intermittent connectivity and remote recovery paths change response assumptions.\\
Multi-cloud & Correlate identities and timestamps across multiple control planes.\\
\bottomrule\end{tabularx}\par\medskip
The handbook addresses these environments as distinct technical subjects rather than a single universal security stack. \handbookref{18--23,83--93,98,104}

\section{On-call shift checklist and post-incident worksheet}
\textbf{Shift start:} collector health; priority open cases; known maintenance; threat advisories; expired allowlists; cloud identity anomalies; backups. \textbf{Shift end:} handoff owner, case status, unresolved hypotheses, evidence location, impact, next action and deadline. \designed{}
\begin{sourcebox}\textbf{Case record fields:} incident ID, detection ID, classification, time source/time zone, investigator, evidence IDs and hashes, scoped entities, containment approvals, recovery verification, communications log, causal findings and owners. \designed{}\end{sourcebox}

\section{Practical verification exercises}
\begin{enumerate}
\item Given a sample five-tuple flow log, identify blocked vs observed flows without assuming packet payloads are known.
\item Design a detector for unexpected privileged cloud API calls; provide three benign counterexamples.
\item Build a coverage map for credential theft and lateral movement, stating where telemetry is absent.
\item Define a controlled restore exercise for an unavailable logging service.
\end{enumerate}
Rubric: evidence fidelity (30), response safety (25), test coverage (25), operational ownership (20). \designed{}

\section*{Handbook chapter alignment}\addcontentsline{toc}{section}{Handbook chapter alignment}
These chapters ground the domain scope. The document offers original synthesis and technical procedures, not a chapter-by-chapter reproduction.
\par\medskip\noindent\begin{longtable}{@{}p{0.085\textwidth}p{0.68\textwidth}p{0.16\textwidth}@{}}\toprule Chapter & Handbook topic & Combined PDF p. \\ \midrule\endhead5 & Detecting System Intrusions & 91 \\
6 & Intrusion Detection in Contemporary Environments & 135 \\
7 & Preventing System Intrusions & 157 \\
8 & Guarding Against Network Intrusions & 175 \\
9 & Fault Tolerance and Resilience in Cloud Computing Environments & 191 \\
11 & UNIX and Linux Security & 233 \\
12 & Eliminating the Security Weakness of Linux and UNIX Operating Systems & 253 \\
13 & Internet Security & 271 \\
14 & The Botnet Problem & 297 \\
15 & Intranet Security & 309 \\
16 & Local Area Network Security & 331 \\
17 & Wireless Network Security & 351 \\
18 & Wireless Sensor Network Security: The Internet of Things & 367 \\
19 & Security for the Internet of Things & 389 \\
20 & Cellular Network Security & 405 \\
21 & Radio Frequency Identification Security & 425 \\
22 & Optical Network Security & 443 \\
23 & Optical Wireless Security & 451 \\
42 & Cyber Forensics and Incidence Response & 717 \\
44 & Network Forensics & 769 \\
59 & Virtual Private Networks & 1065 \\
61 & VoIP Security & 1093 \\
64 & Securing Cloud Computing Systems & 1151 \\
65 & Cloud Security & 1167 \\
66 & Cloud Security Trends & 1179 \\
67 & Private Cloud Security & 1195 \\
68 & Virtual Private Cloud Security & 1203 \\
69 & Protecting Virtual Infrastructure & 1211 \\
70 & Software-Defined Networking and Network Function Virtualization Security & 1219 \\
74 & Intrusion Prevention and Detection Systems & 1289 \\
75 & Transmission Control Protocol/Internet Protocol Packet Analysis & 1303 \\
76 & Firewalls & 1317 \\
78 & System Security & 1345 \\
80 & Endpoint Security & 1361 \\
83 & Securing the Infrastructure & 1387 \\
84 & Cyberattacks Against the Grid Infrastructure & 1401 \\
85 & Energy Infrastructure Cyber Security & 1409 \\
86 & Homeland Security and Cybersecurity & 1417 \\
87 & Cyber Warfare & 1431 \\
88 & Cyber Attack Process & 1467 \\
89 & Smart Cities: Cybersecurity Concerns & 1483 \\
90 & Disaster Management Cycle, Smart Technologies, Urban Form, Traditional Developments, and Sprawl & 1499 \\
91 & Cyber Security in Smart Homes & 1523 \\
92 & Future Trends for Cyber Security for Smart Cities and Homes & 1543 \\
93 & An Overview of Cyber Attacks and Defenses on Intelligent Connected Vehicles & 1567 \\
98 & Satellite Cyber Attack Search and Destroy & 1647 \\
103 & Future Trends for Cyber Attacks in the Healthcare Industry & 1737 \\
104 & Future Trends in Maritime Cybersecurity & 1749 \\
\bottomrule\end{longtable}\par\medskip

\section*{Source and reference register}
\addcontentsline{toc}{section}{Source and reference register}
\begin{description}[leftmargin=1.6em,style=nextline]
\item[H] John R. Vacca (ed.), \emph{Computer and Information Security Handbook}, 4th ed., Volumes I--II, Elsevier/Morgan Kaufmann, 2025. Chapter identifiers in this document point to this book.
\item[N1] NIST, \emph{Cybersecurity Framework 2.0} (2024), \url{https://www.nist.gov/cyberframework}.
\item[N2] NIST SP 800-218, \emph{Secure Software Development Framework v1.1} (2022), \url{https://csrc.nist.gov/pubs/sp/800/218/final}. Later drafts are not treated as final requirements.
\item[N3] NIST SP 800-61 Rev. 3, \emph{Incident Response Recommendations and Considerations for Cybersecurity Risk Management} (2025), \url{https://csrc.nist.gov/pubs/sp/800/61/r3/final}.
\item[N4] NIST SP 800-63-4, \emph{Digital Identity Guidelines} (2025), \url{https://csrc.nist.gov/pubs/sp/800/63/4/final}.
\item[N5] NIST, FIPS 203/204/205 (2024), \url{https://csrc.nist.gov/Projects/Post-Quantum-Cryptography/publications}.
\item[O1] OWASP, \emph{API Security Top 10} (2023), \url{https://api-security.owasp.org/editions/2023/en/0x11-t10/}.
\item[C1] CISA, \emph{Known Exploited Vulnerabilities Catalog}, \url{https://www.cisa.gov/known-exploited-vulnerabilities-catalog}.
\item[M1] MITRE ATT\&CK Enterprise, \url{https://attack.mitre.org/matrices/enterprise/}.
\item[G1] GitHub, \emph{REST API endpoints for code scanning}, \url{https://docs.github.com/en/rest/code-scanning}.
\end{description}
\begin{warningbox}\textbf{Interpretation rule.} \handbookref{number} identifies a conceptual anchor in the supplied handbook, not a verbatim quotation or an official requirement. \modernext{} identifies a contemporary extension informed by external references or current engineering practice. \designed{} identifies an independently developed example, template, exercise, metric, or operational procedure. Requirements are binding only if separately adopted through applicable law, contract, policy, or the cited standards; this suite does not assert otherwise.\end{warningbox}

\end{document}
